%% 12장 수와 글자 사이
\chapter{수와 글자 사이}
\label{ch:number-text}

\abstract{데이터는 대개 글자로 들어와 글자로 나간다. CSV 파일의 \code{0.1}이라는 세 글자는 binary64 수로 바뀌어야 계산에 쓰이고, 계산 결과는 다시 글자로 바뀌어 화면과 파일에 적힌다. 이 장은 그 두 방향의 변환을 다룬다. 십진수를 binary64로 바꾸는 일은 가장 가까운 binary64 수를 찾는 문제이고, 이진수를 십진수로 바꾸는 일은 다시 읽었을 때 같은 수가 되는 가장 짧은 십진 표현을 찾는 문제이다. 두 문제 모두 정의는 단순하지만 빠르면서 언제나 정확하게 풀기는 어렵다. Clinger의 빠른 길과 Eisel--Lemire의 알고리즘으로 읽기를, 가장 짧은 왕복 표현의 성질과 Ryū의 생각으로 쓰기를 살핀 뒤, 여덟 자리를 한꺼번에 읽는 비트 기교를 보고, 마지막으로 글자 자체를 바이트로 적는 UTF-8 부호화와 그 검증을 다룬다.}

\section{십진수와 이진수}
\label{sec:decimal-binary}

십진수 $0.1$이 binary64로 정확히 나타나지 않는 이유는 간단하다. 기약분수 $m/n$이 유한한 이진 소수로 적히는 것은 $n$이 $2$의 거듭제곱일 때뿐인데, $1/10$의 분모에는 인수 $5$가 있다. 거꾸로 모든 binary64 수는 분모가 $2$의 거듭제곱이고 $10=2\cdot5$이므로 유한한 십진 소수로 정확히 적힌다. 8장의 처음에 본 $0.1$의 binary64 값 $0.1000000000000000055511\cdots$은 소수점 아래 $55$자리에서 끝난다. 그러니 이진수에서 십진수로 가는 변환은 언제나 정확하게 \textbf{할 수는} 있지만 그대로 적으면 너무 길고, 십진수에서 이진수로 가는 변환은 대개 반올림을 피할 수 없다.

두 변환에 요구하는 성질을 분명히 해 두자. 십진수 $d$를 읽을 때는 $d$에 가장 가까운 binary64 수 $\fl(d)$를 내야 하는데 이것을 \term{올바른 반올림 읽기}{correctly rounded parsing}라 한다. binary64 수 $x$를 쓸 때는 다시 읽었을 때 $x$가 되는, 곧 $\fl(d)=x$인 십진수 $d$ 가운데 자릿수가 가장 적은 것을 낸다. 이것이 \term{가장 짧은 왕복 표현}{shortest round-trip representation}이다. 두 성질을 함께 지키면 수를 글자로 썼다가 다시 읽어도 비트 하나 바뀌지 않으며, Python의 \code{repr}과 \code{float}가 바로 이 약속을 지킨다.

\subsection{몇 자리면 충분한가}

binary64 수를 십진 $k$자리로 적었다가 다시 읽을 때 언제나 같은 수가 되려면 $k$가 얼마여야 할까. 십진 $k$자리 수들의 간격이 binary64 수들의 간격보다 충분히 좁으면 된다. 이 조건을 정리하면 다음을 얻는다\cite{matula1968}.

\begin{theorem}[Matula]
\label{thm:matula}
모든 binary64 수는 유효숫자 $17$자리의 십진수로 적었다가 올바른 반올림으로 다시 읽으면 원래의 수로 돌아온다. 거꾸로 유효숫자 $15$자리 이하의 십진수는 binary64로 올바르게 반올림했다가 같은 자릿수의 십진수로 올바르게 반올림하면 원래의 십진수로 돌아온다. 두 수 $17$과 $15$는 모두 더 줄일 수 없다.
\end{theorem}

증명의 요점은 간격의 비교이다. $p$비트 이진수의 상대 간격은 $2^{-p}$에서 $2^{1-p}$ 사이이고 $k$자리 십진수의 상대 간격은 $10^{-k}$에서 $10^{1-k}$ 사이이므로, $10^{1-k}<2^{-p}$, 곧 $k>1+p\log_{10}2$이면 십진 쪽이 어디서나 더 촘촘해서 두 binary64 수 사이에 늘 십진수가 하나 이상 끼어든다. $p=53$이면 $1+53\log_{10}2\approx16.95$이므로 $k=17$이다. 거꾸로 방향은 $2^{1-p}<10^{-k}$, 곧 $k<(p-1)\log_{10}2\approx15.65$에서 $k=15$를 얻는다. 실제로 $0.1+0.2$의 binary64 값은 $15$자리로 쓰면 \code{0.3}, $16$자리로 써도 \code{0.3000000000000000}이어서 다시 읽으면 다른 수가 되고, $17$자리 \code{0.30000000000000004}가 되어서야 원래 수로 돌아온다.

$17$자리면 언제나 충분하지만 대부분의 수에는 더 적은 자리로 충분하다. \code{0.1}의 binary64 값은 $17$자리로 쓰면 $0.10000000000000001$이지만 한 자리 \code{0.1}만으로도 다시 읽으면 같은 수가 된다. 그래서 가장 짧은 왕복 표현이 사람에게 보여 주기에 좋다. 반대로 binary64 수가 아닌 십진수를 적었을 때는 그 십진수가 그대로 나오지 않을 수 있다. $10^{23}$은 binary64 수가 아니고 그에 가장 가까운 수는 $10^{23}-8388608$인데, 이 수의 가장 짧은 왕복 표현은 \code{1e+23}이다. 겉보기에는 정확한 것 같지만 이 글자가 가리키는 binary64 수는 $10^{23}$과 다르다.

\section{십진수 읽기}
\label{sec:parsing}

\subsection{Clinger의 빠른 길}

십진수 텍스트를 $w\times10^q$로 읽었다고 하자. 여기서 $w$는 유효숫자를 모은 정수이고 $q$는 십진 지수이다. 예컨대 \code{12.5e3}은 $w=125$, $q=2$이다. $w$와 $10^{|q|}$이 둘 다 binary64 수이면, $w\times10^q$는 binary64 연산 하나($q\ge0$이면 곱셈 하나, $q<0$이면 나눗셈 하나)이고, 8장에서 보았듯이 그 연산은 올바르게 반올림된다. 따라서 결과가 곧 $\fl(w\times10^q)$이다. $w$는 $2^{53}$ 이하이면 정확하고, $10^q=2^q5^q$은 $5^q<2^{53}$이면 정확한데 $5^{22}=2384185791015625<2^{53}$이고 $5^{23}>2^{53}$이므로 $|q|\le22$이다. 이 조건에서 연산 하나로 끝내는 것이 \term{Clinger의 빠른 길}{Clinger's fast path}이다\cite{clinger1990}. CSV 파일의 수 대부분은 유효숫자가 많지 않고 지수가 작으므로 이 길로 충분하다.

빠른 길을 벗어나면 일반적으로는 큰 정수 산술이 필요하다. $w\times10^q$를 정확한 유리수로 두고 가장 가까운 binary64 수를 찾으면 되지만, $q$가 수백이면 수백 자리의 정수를 다루어야 한다. Clinger와 Gay는 근삿값을 먼저 구하고 그것이 맞는지를 큰 정수 비교로 확인하는 방법을 정리했는데\cite{clinger1990,gay1990}, 이 방법은 정확하지만 느리다.

\subsection{Eisel--Lemire 알고리즘}

Eisel과 Lemire는 큰 정수 없이 거의 모든 경우를 $128$비트 곱셈 하나로 끝내는 방법을 찾았다\cite{lemire2021parsing}. 생각은 이렇다. $10^q=2^q\cdot5^q$에서 $2^q$는 지수에 더하면 되므로, $w\cdot5^q$의 앞쪽 $54$비트를 아는 일이 어렵다. $5^q$의 앞쪽 $128$비트를 미리 표로 계산해 두고, $w$를 $64$비트 정수의 맨 위 비트가 $1$이 되도록 왼쪽으로 밀어 둔 뒤 둘을 곱하면, 곱의 위쪽 $64$비트가 $w\cdot5^q$의 유효숫자의 앞부분이다. binary64의 유효숫자는 $53$비트이고 반올림에 한 비트가 더 필요하므로 이 $64$비트에서 $10$비트 남짓이 남는다. 표의 $5^q$가 잘린 값이라 곱도 참값보다 조금 작을 수 있지만, 그 오차는 곱의 아래쪽 비트에만 영향을 준다. 그래서 반올림에 쓰는 비트 바로 아래의 남는 비트들이 모두 $1$인, 곧 참값이 반올림 경계에 아주 가까울 수 있는 드문 경우만 아니면 이 곱만으로 반올림이 결정된다. 드문 경우에는 표의 다음 $64$비트까지 곱해 다시 확인하고, 그래도 결정되지 않으면 느린 일반 알고리즘으로 넘긴다.

이진 지수도 곱셈 하나로 구한다. $10^q$의 이진 지수는 $\lfloor q\log_210\rfloor$인데, $\log_210\cdot2^{16}\approx217705.88$이므로 쓰이는 범위의 $q$ 전체에서 정수 $217706$을 곱하고 $16$비트 오른쪽으로 미는 고정소수점 계산이 정확한 값을 준다. 반올림의 마지막 단계는 짝수로의 반올림을 그대로 흉내 내야 하는데, 정확히 한가운데인 경우는 $q$가 작은 범위에서만 생길 수 있으므로 그 범위에서만 따로 확인한다.

\begin{instardust}[\code{csrc/frame/ops/parse.hpp}]
데이터프레임의 \code{read\_csv}와 문자열 열의 \code{astype(float64)}는 같은 파서를 쓰고, 파일의 주석은 fast\_float의 알고리즘을 논문에서 직접 옮겨 적었다고 밝힌다. \code{eisel\_lemire} 함수는 위의 단계 그대로이다. $w$를 맨 위 비트까지 밀고, \code{pow5\_table.hpp}의 $5^q$ $128$비트 근삿값과 곱하고, 위쪽 $64$비트의 아래 $9$비트가 모두 $1$이면 표의 다음 단어를 더 곱한다. 이진 지수는 $(152170+65536)\cdot q$를 $16$비트 민 값에 $63$을 더해 구하는데, $152170+65536=217706$이다. 설계 문서는 이 파서가 무작위 텍스트 $2000$만 개에서 Python의 \code{float(text)}와 비트까지 같았다고 기록한다. 유효숫자가 $19$자리를 넘거나 곱이 반올림을 결정하지 못하는 드문 경우는 C++ 표준 라이브러리의 \code{from\_chars}로 넘긴다.
\end{instardust}

\subsection{여덟 자리를 한꺼번에}

숫자 문자열을 정수로 바꾸는 가장 단순한 방법은 한 글자씩 $w\gets10w+(c-\code{'0'})$을 되풀이하는 것이다. 이 반복은 매 단계가 앞 단계를 기다리므로 느리고, 글자마다 숫자인지 확인하는 분기가 붙는다. 64비트 레지스터 하나에 ASCII 글자 여덟 개가 들어간다는 사실을 쓰면 이것을 크게 줄일 수 있는데, 하나의 레지스터 안에서 여러 개의 작은 값을 동시에 다루는 이런 기법을 \term{레지스터 안의 SIMD}{SWAR, SIMD within a register}라 한다.

\code{"12345678"}의 여덟 바이트를 작은 쪽이 먼저 오는 순서로 읽으면 맨 아래 바이트가 \code{'1'}인 64비트 정수가 된다. 바이트마다 \code{'0'}$=\mathtt{0x30}$을 빼면 각 바이트가 숫자값 하나를 담는다.
\[
\mathtt{0x0807060504030201}.
\]
여기에 $10$을 곱하고 $8$비트 오른쪽으로 민 자신을 더하면, 이웃한 두 바이트가 $10a+b$로 합쳐져 짝수 번째 바이트마다 두 자리 수 $12,34,56,78$이 생긴다(홀수 번째 바이트에는 쓰지 않는 값이 남는다). 같은 방식으로 이웃한 두 자리 수를 $100a+b$로, 다시 네 자리 수를 $10000a+b$로 합치는데, Lemire의 구현은 둘째와 셋째 단계를 상수 두 개와의 곱셈 두 번과 $32$비트 밀기로 한꺼번에 해 $12345678$을 얻는다. 여덟 바이트가 모두 숫자인지도 분기 없이 확인할 수 있다. 각 바이트에 $\mathtt{0x46}$을 더한 값과 $\mathtt{0x30}$을 뺀 값의 논리합에서 바이트마다 맨 위 비트가 하나라도 서면 그 바이트는 \code{'0'}부터 \code{'9'}까지의 범위 밖에 있다. 여덟 글자를 처리하는 데 산술 연산 열 개 남짓이면 되고 분기는 없다.

\section{이진수 쓰기}
\label{sec:printing}

binary64 수 $x$의 가장 짧은 왕복 표현을 찾는 문제를 생각하자. $x$의 이웃한 두 binary64 수를 $x^-$와 $x^+$라 하면, 다시 읽어서 $x$가 되는 실수는 $x$와 두 이웃의 중점 사이, 곧 반올림 구간
\[
\Bigl[\frac{x^-+x}2,\ \frac{x+x^+}2\Bigr]
\]
에 있는 수들이다(양 끝이 포함되는지는 짝수로의 반올림 규칙이 정한다). 따라서 가장 짧은 왕복 표현은 이 구간 안에 있는 십진수 가운데 유효숫자가 가장 적은 것이고, 그런 것이 여럿이면 $x$에 가장 가까운 것을 고른다. 이 구간의 양 끝을 정확히 다루는 데서 문제가 생긴다. 구간의 끝은 binary64 수가 아니고 그 십진 표현은 아주 길 수 있기 때문이다.

Steele과 White는 이 문제를 큰 정수로 정확하게 푸는 Dragon4 알고리즘을 1990년에 발표했고\cite{steele1990}, 이후 수십 년 동안 그것을 빠르게 하는 연구가 이어졌다. Adams의 Ryū는 이 구간을 $10$의 거듭제곱으로 크기를 맞춘 정수 구간으로 바꾼 뒤, 큰 정수 대신 미리 계산한 $5^q$의 근삿값과의 $128$비트 곱셈으로 구간의 끝을 충분히 정확하게 계산해서, 구간에 들어가는 가장 짧은 십진수를 몇 번의 나눗셈만으로 찾는다\cite{adams2018ryu}. Eisel--Lemire와 같은 종류의 생각으로, $5^q$의 표와 $128$비트 곱으로 큰 정수를 대신하되 그것으로 결정되지 않는 경우가 없거나 드물다는 것을 증명하는 방식이다. 오늘날 C++ 표준 라이브러리의 \code{to\_chars}와 많은 언어의 수 출력이 Ryū나 그 후속 알고리즘을 쓴다.

\begin{instardust}[\code{csrc/frame/ops/cast.cpp}]
데이터프레임의 float 열을 문자열로 바꾸는 \code{python\_float\_str}은 Python의 \code{repr(float)}과 같은 글자를 낸다. 가장 짧은 왕복 숫자열은 \code{std::to\_chars}의 지수 꼴에서 얻고, 소수점의 위치를 따져 십진 지수가 $(-4,16]$ 안이면 고정 소수점으로, 밖이면 \code{d.ddde+XX} 꼴로 다시 적으며, 정수인 고정 소수점 값에는 \code{.0}을 붙인다. NaN과 무한대는 \code{nan}, \code{inf}, \code{-inf}이다. 읽기 쪽이 Python의 \code{float}와 비트까지 같고 쓰기 쪽이 \code{repr}과 글자까지 같으므로, stardust로 쓴 CSV를 Python으로 읽거나 그 반대로 해도 수가 바뀌지 않는다.
\end{instardust}

\section{글자를 바이트로}
\label{sec:utf8}

\subsection{UTF-8}

컴퓨터의 글자는 \term{코드 포인트}{code point}라 부르는 $0$부터 $\mathtt{0x10FFFF}$까지의 정수이고, 이것을 바이트열로 적는 규칙이 부호화이다. Unicode는 세상의 거의 모든 문자에 코드 포인트를 부여하는데\cite{unicode17}, 오늘날 파일과 네트워크에서 가장 널리 쓰이는 부호화가 \term{UTF-8}{UTF-8}이다\cite{rfc3629}. UTF-8은 코드 포인트의 크기에 따라 $1$바이트에서 $4$바이트를 쓴다.

\begin{table}[t]
\begin{sdtable}{colspec={lll}, row{1}={font=\small\bfseries}}
\toprule
코드 포인트의 범위 & 바이트 꼴 & 예 \\
\midrule
\texttt{U+0000}--\texttt{U+007F} & \texttt{0xxxxxxx} & \texttt{A} $\to$ \texttt{41} \\
\texttt{U+0080}--\texttt{U+07FF} & \texttt{110xxxxx 10xxxxxx} & \texttt{é} $\to$ \texttt{C3 A9} \\
\texttt{U+0800}--\texttt{U+FFFF} & \texttt{1110xxxx 10xxxxxx 10xxxxxx} & \texttt{한} $\to$ \texttt{ED 95 9C} \\
\texttt{U+10000}--\texttt{U+10FFFF} & \texttt{11110xxx 10xxxxxx}$\times3$ & \texttt{U+1F600} $\to$ \texttt{F0 9F 98 80} \\
\bottomrule
\end{sdtable}
\caption{UTF-8의 바이트 꼴. \texttt{x} 자리에 코드 포인트의 비트를 위에서부터 채운다. 한글 음절 `한'의 코드 포인트는 \texttt{U+D55C}이다.}
\label{tab:utf8}
\end{table}

\cref{tab:utf8}의 설계에는 몇 가지 영리한 성질이 있다. ASCII 글자는 원래의 한 바이트 그대로이므로 ASCII 파일은 그 자체로 올바른 UTF-8이다. 여러 바이트로 된 글자의 바이트는 모두 맨 위 비트가 $1$이므로, ASCII 글자를 찾는 검색, 예컨대 CSV에서 쉼표와 줄바꿈을 찾는 일이 다른 글자의 중간 바이트에 속아 넘어가지 않는다. 첫 바이트는 \texttt{0}, \texttt{110}, \texttt{1110}, \texttt{11110}으로 시작하고 이어지는 바이트는 언제나 \texttt{10}으로 시작하므로, 바이트열의 아무 곳에서 읽기 시작해도 글자의 경계를 곧바로 찾을 수 있다. 또 바이트열을 바이트 단위로 사전식 비교하면 코드 포인트의 순서와 같은 순서가 나온다.

\subsection{검증}

모든 바이트열이 올바른 UTF-8은 아니다. 이어지는 바이트가 \texttt{10}으로 시작하지 않거나 바이트가 모자라는 경우는 분명한 오류이고, 그 밖에 세 가지가 더 금지된다. 첫째는 \term{과잉 부호화}{overlong encoding}로, 예컨대 \texttt{/}($\mathtt{0x2F}$)를 두 바이트 \texttt{C0 AF}로 적는 것처럼 더 짧게 적을 수 있는 코드 포인트를 길게 적은 것이다. 이것을 허용하면 같은 글자가 여러 바이트열로 적힐 수 있어서, 특정 글자를 걸러 내는 보안 검사를 우회하는 데 악용된 적이 있다. 둘째는 UTF-16에서 쓰는 대리 영역 \texttt{U+D800}--\texttt{U+DFFF}이고, 셋째는 \texttt{U+10FFFF}를 넘는 값이다\cite{rfc3629}.

검증의 가장 큰 비용은 대부분이 ASCII인 텍스트를 한 바이트씩 확인하는 데서 나온다. 그래서 여덟 바이트를 64비트 정수 하나로 읽어 $\mathtt{0x8080808080808080}$과 논리곱해 보는 것만으로 여덟 바이트가 모두 ASCII인지를 한 번에 확인하고, ASCII가 아닌 바이트를 만난 곳부터만 글자 단위로 꼼꼼히 확인하는 것이 흔한 방법이다.

\begin{instardust}[\code{csrc/frame/io/csv.cpp}]
\code{valid\_utf8}은 바로 이 두 단계로 되어 있다. 여덟 바이트씩 ASCII인지 확인하다가 맨 위 비트가 선 바이트를 만나면, 첫 바이트로 길이를 정하고 이어지는 바이트가 \texttt{10}으로 시작하는지 본 뒤 코드 포인트를 모아 과잉 부호화, 대리 영역, \texttt{U+10FFFF} 초과를 거부한다. \code{read\_csv}는 한 걸음 더 나아가 파일을 읽으면서 $\mathtt{0x80}$ 이상의 바이트가 하나도 없었는지를 기억해 두고, 그런 파일에서는 문자열 열의 검증을 아예 건너뛴다. 파일 맨 앞의 UTF-8 바이트 순서 표시 \texttt{EF BB BF}는 내용이 아니므로 떼어 내며, 2장에서 본 따옴표의 홀짝을 계산하는 비트 기교도 이 표시 안에서는 따옴표를 세지 않도록 맞추어 둔다.
\end{instardust}

\begin{aside}[식탁 위의 UTF-8]
UTF-8은 1992년 9월 Ken Thompson과 Rob Pike가 Bell 연구소의 Plan 9 운영체제를 위해 고안했다. 전해지는 이야기로는 X/Open의 국제화 회의에서 제안된 부호화가 마음에 들지 않았던 두 사람이 저녁 식사 자리에서 식탁 매트 위에 지금의 설계를 그렸고, 며칠 만에 Plan 9 전체를 그 부호화로 바꾸었다고 한다. ASCII와의 호환, 자기 동기화, 바이트 순서와 무관함이라는 이 장에서 본 성질들이 그날 저녁의 설계에 이미 모두 들어 있었고, 그 덕분에 UTF-8은 오늘날 웹 문서의 대부분을 차지하는 부호화가 되었다.
\end{aside}

\section{맺음}

수와 글자 사이의 변환은 데이터 분석의 가장 앞과 가장 뒤에 있어서 눈에 잘 띄지 않지만, 여기서 비트 하나가 어긋나면 그 뒤의 모든 계산이 다른 답을 낸다. 올바른 반올림 읽기와 가장 짧은 왕복 쓰기라는 두 약속은 이 위험을 없애고, Eisel--Lemire와 Ryū 같은 알고리즘은 그 약속을 지키면서도 변환을 데이터를 읽는 속도의 병목에서 빼냈다. 표와 $128$비트 곱으로 큰 정수를 대신하는 생각과, 레지스터 하나에 여러 값을 담아 한꺼번에 처리하는 생각은 다음 장의 컴퓨터 구조에서 다시 나온다.

\chapterreferences
